\chapter{SURROGATE MODELLING AND OPTIMIZATION}
\label{ch:method}

We parameterize a laminate by its ply angles and fibre volume fraction, train a Gaussian process
on finite-element results, and use expected improvement to propose layups that minimize mass at a
required stiffness.

\begin{figure}[htbp]
  \centering
  \begin{tikzpicture}[node distance=5mm and 7mm,
      box/.style={draw=accent, rounded corners=2pt, minimum height=9mm, align=center, font=\small}]
    \node[box] (d) {Layup\\parameters};
    \node[box, right=of d] (f) {Finite-element\\simulation};
    \node[box, right=of f] (g) {Gaussian-process\\surrogate};
    \node[box, right=of g] (b) {Bayesian\\optimization};
    \draw[-{Stealth}] (d) -- (f); \draw[-{Stealth}] (f) -- (g); \draw[-{Stealth}] (g) -- (b);
  \end{tikzpicture}
  \caption{Surrogate-assisted design loop.}
  \label{fig:loop}
\end{figure}
